% 本文件由 assemble.py 按 _work/plans/ch90.json 逐字组装，请勿手改（改 plan 或 blocks 后重新生成）。
\chapter{数学与物理方法（通用）}
\label{ch:90}

\section{常用数学工具}
% ---- block 091-08 (原稿第91页) kind=formula ----
\textbf{二项式定理}\quad $(1+x)^{\alpha}\approx 1+\alpha x$\quad ($x$ 很小时)

% ---- block 199-06 (原稿第199页) kind=formula ----
$(\tan x)'=\dfrac{1}{\cos^2 x}$

\notefig[0.55]{figures/p199_e.png}
% 图为 $y=\tan x$ 的导函数 $\frac{1}{\cos^2x}$ 的草图（竖直虚线为渐近线），箭头自公式指向中间一支；右侧另有一小段红笔画的 V 形曲线

% ---- block 001-01 (原稿第1页) kind=derivation ----
\[
v=kx,\quad \frac{\dd x}{\dd t}=kx,\quad \int\frac{\dd x}{kx}=\int\dd t,\quad \ln x=kt+kc,\quad x=\ee^{kt+kc}
\]

% ---- block 091-06 (原稿第91页) kind=method ----
\textbf{数列类}
\[
\left\{\begin{aligned}
S_n&=\frac{(a_1+a_n)n}{2}\\
S_n&=\frac{a_1}{q-1}(q^{n}-1)\qquad \text{若}\ |q|<1\ \ \lim_{n\to\infty}S_n=\frac{a_1}{1-q}
\end{aligned}\right.
\]
\begin{quote}
详写过程 $|$ 易得过程2\\
综上全过程得出通项式\\
求前$n$项和、极限$n$无穷
\end{quote}
% CHECK: “过程”后的竖线也可能是数字 1（详写过程1）


\section{物理量的定义式与变化率}
% ---- block 083-07 (原稿第83页) kind=summary ----
% 原稿左侧有大括号：{定义式 比值 (无关)；决定式 (相关)}，“无关”“相关”均被圈起；公式按原稿行序排列
\[
\left\{\begin{aligned}
&\text{定义式\ \ 比值}\\[8pt]
&\text{决定式}
\end{aligned}\right.
\]

$v=\dfrac{\Delta x}{\Delta t}$\quad $k=\dfrac{F}{x}$\quad $E=\dfrac{F}{q}$\quad $\unclear{R=\dfrac{U}{I}}$ % CHECK: 末式笔迹形似“k=v/L”，按R=U/I转录(本页“R”常写得像“k”)

$I=\dfrac{q}{t}$\quad $C=\dfrac{Q}{U}$\quad $\Phi=B_{\perp}S$

\fbox{无关}\quad $p=mv$\quad $I=Ft$

\fbox{相关}\quad $a=\dfrac{F}{m}$\quad $E=\dfrac{kQ}{r^2}$

$I=nesv$\quad $R=\dfrac{\rho L}{S}$\quad $C=\dfrac{\varepsilon S}{4\pi kd}$

% ---- block 083-08 (原稿第83页) kind=formula ----
\textbf{函数分析}

% 原稿左侧红笔大括号括住下列各式
\[
\left\{\begin{aligned}
a&=\frac{\Delta v}{\Delta t}\\[6pt]
E_{\text{\unclear{感}}}&=-n\frac{\Delta\Phi}{\Delta t}\qquad E_{\text{场}}=-\red{\frac{\Delta\varphi}{\Delta x}}\\[6pt]
i&=\frac{\Delta q}{\Delta t}\\[6pt]
F&=\frac{\Delta p}{\Delta t}\\[6pt]
P&=\frac{\Delta W}{\Delta t}\\[6pt]
F_{\text{合}}&=\frac{\Delta E_k}{\Delta x}\\[6pt]
F_{\text{其他}}&=\frac{\Delta E_{\text{机}}}{\Delta x}\\[6pt]
R&=\frac{U_1}{I_1}=\frac{U_2}{I_2}\neq\frac{U_2-U_1}{I_2-I_1}\\[6pt]
P&=U_1I_1\qquad\text{面积}\ S=\int xy
\end{aligned}\right.
\]
% 原稿E_场式中黑笔旧写的分式被红笔涂划，右侧红笔重写为Δφ/Δx；E_感下标笔迹似“源/感”

$\dfrac{\Delta v}{\Delta t}$\quad 匀圆变化\quad 方向变


\section{常用解题方法}
% ---- block 221-03 (原稿第221页) kind=method ----
\textbf{3、}代值、极限法

% ---- block 270-08 (原稿第270页) kind=method ----
\textbf{特殊值法}\quad 有$k$、\unclear{$n$}、类参数、比值、直接赋值。

% ---- block 221-10 (原稿第221页) kind=method ----
\textbf{补偿法}
\notefig[0.25]{figures/p221_f.png}
\unclear{$F_{\text{外大小}}=F_{\text{大小}}-F_{\text{小小}}$} % CHECK: 各下标笔迹难辨

\textbf{等效法}
\notefig[0.25]{figures/p221_g.png}

\textbf{图像法}


\section{解题思路与注意事项}
% ---- block 239-03 (原稿第239页) kind=method ----
模型组合，分步解决

平抛重分解，圆周重受力，天体重轨道，功能重对应

% ---- block 163-07 (原稿第163页) kind=method ----
能否到达\ \ 假设到达推矛盾\qquad 相交相遇问题轨迹方程法

% ---- block 267-02 (原稿第267页) kind=note ----
单位\unclear{$\ne$}物理量\qquad 求比值要求绝对值

